
\documentclass[11pt,a4paper]{article}
\usepackage[utf8]{inputenc}
\usepackage[T1]{fontenc}
\usepackage[spanish]{babel}
\usepackage{geometry}
\usepackage{graphicx}
\usepackage{hyperref}
\usepackage{xcolor}
\usepackage{tabularx}
\geometry{margin=2cm}
\setlength{\parindent}{0pt}
\setlength{\parskip}{0.55em}
\sloppy
\definecolor{tcTeal}{HTML}{0E7C74}
\definecolor{tcDark}{HTML}{101817}
\definecolor{tcMuted}{HTML}{526A65}
\definecolor{tcWarnBg}{HTML}{FFF7E5}
\definecolor{tcWarnBorder}{HTML}{C99123}
\hypersetup{colorlinks=true, urlcolor=tcTeal, linkcolor=tcTeal}
\begin{document}

\begin{center}
{\huge\bfseries Informe AI Analyst: US30 H1 rsi extreme with context}\\[0.25em]
{\large Reporte read-only para revision humana}
\end{center}

\vspace{0.2em}
\hrule height 0.7pt
\vspace{0.9em}

\begin{tabularx}{\textwidth}{>{\bfseries}p{3.2cm}X}
Paquete & US30 H1 rsi extreme with context 86844a61 \\
Setup & US30 H1 rsi extreme with context \\
Modelo & gpt-5.5 \\
Thinking & medium \\
Macro & high \\
Estado & reviewed \\
\end{tabularx}

\vspace{0.75em}
\fcolorbox{tcWarnBorder}{tcWarnBg}{%
\begin{minipage}{0.94\textwidth}
\textbf{Limite de uso.} Este informe no es una senal, no autoriza MT5 y no ejecuta ordenes. Sirve para priorizar revision humana y documentar el contexto.
\end{minipage}}

\vspace{0.4em}

\section*{Resumen}
Revision read-only de US30 en H1. El paquete muestra un contexto de sobreventa en RSI H1 dentro de una tendencia mayor H1/H4/D1 aun alcista, pero el tramo inmediato viene de una caida brusca desde zona de maximos previos y no aporta timing operativo propio. Estado adecuado: observacion y revision humana, sin ejecucion ni aprobacion operativa.

\section*{Lectura tecnica}
El precio cierra en 50.998,45 tras un retroceso intenso desde el maximo intradia de 51.835,8. La lectura principal es de pullback profundo dentro de una estructura mayor todavia constructiva: H1, H4 y D1 figuran alcistas, mientras M15 ya esta bajista, lo que introduce friccion de corto plazo. El RSI H1 esta en 23,77 y el estocastico H1 cerca de 10, lo que confirma agotamiento bajista de corto plazo, pero el MACD H1 se ha girado con fuerza a la baja y visualmente aun no muestra recuperacion clara.



\section*{Grafico de referencia}
\begin{center}
\includegraphics[width=0.96\textwidth]{\detokenize{ai_analyst_report_chart.png}}
\end{center}


\section*{Lectura de confluencias}
RSI H1 en zona de sobreventa, coherente con la etiqueta de extremo alcista para revision; Tendencia H1/H4/D1 alineada al alza, lo que da contexto de pullback y no de giro mayor confirmado; Precio muy cerca de la zona Fibonacci 78,6 y de los niveles redondos 50.900-51.000; Volatilidad H1 adecuada, con ratio 1,10 frente a su mediana, suficiente para que la lectura sea relevante sin ser un desorden extremo por si sola; y Alta correlacion estructural con US500 en D1/H4/H1, util para contrastar si el movimiento es de mercado amplio y no solo del Dow.

\section*{Contradicciones y cautelas}
M15 esta bajista mientras H1/H4/D1 permanecen alcistas; la estructura corta aun no acompana la lectura de recuperacion; El cierre queda por debajo de varios niveles Fibonacci intermedios y cerca del retroceso profundo, mostrando deterioro reciente aunque el marco mayor siga positivo; La lectura de onda candidata marca un posible W2? en H1, pero su direccion del tramo actual es bajista y la calidad es media, por lo que no debe sobreponderarse; El MACD H1 cae con pendiente negativa tras el impulso previo, contradiciendo una lectura de rebote inmediato basada solo en RSI; y El paquete esta en estado watching y sin nivel de disparo; por tanto la calidad contextual no equivale a timing operativo.

\section*{Riesgos principales}
Paquete study-only: no es una senal, no habilita ejecucion y no debe conectarse a MT5 ni Telegram; La vela final de la ventana tiene rango amplio y cierre cerca de minimos relativos, lo que aumenta el riesgo de continuidad bajista antes de cualquier estabilizacion; La zona 50.900-51.000 concentra nivel redondo y Fibonacci cercano; una perdida limpia de esa zona dejaria expuesto el minimo previo y soportes inferiores; El spread y la actividad aumentan durante las horas de mayor movimiento, por lo que la lectura visual debe contrastarse con condiciones reales de mercado si se revisa manualmente; y Las correlaciones preservadas como ultimo dato bueno tienen valor contextual, pero conviene comprobar datos actualizados antes de usar relaciones intermercado en la revision.

\section*{Contexto macro y noticias}
El contexto macro verificado alrededor del cierre del 5 de junio de 2026 es exigente para indices estadounidenses. AP reporto que el Dow Jones Industrial Average cayo 695,15 puntos, hasta 50.866,78, en una sesion marcada por ventas amplias y presion en tecnologia; la explicacion central fue un informe laboral mas fuerte de lo esperado que impulso los rendimientos y endurecio la lectura de politica monetaria. Para US30, esto encaja con la vela bajista final del paquete: la sobreventa tecnica aparece tras un evento macro que pudo aumentar la prima de riesgo y no necesariamente por una simple correccion tecnica aislada.

La lectura de tipos tambien queda mas sensible. Fuentes de calendario economico situan el informe de empleo de Estados Unidos el 5 de junio de 2026, y la cobertura de mercado indico que el dato de nominas reforzo expectativas de una Reserva Federal menos acomodaticia. Ademas, el calendario oficial de la Fed mantiene reuniones de politica monetaria en junio de 2026, por lo que el mercado puede seguir repricing bonos, dolar y renta variable antes de la decision. En una revision humana, este contexto obliga a tratar la confluencia tecnica como vulnerable a titulares y movimientos de rendimiento, sin convertir la sobreventa en permiso operativo.

Desde una lectura financiera, este contexto no cambia por si solo el setup tecnico, pero si condiciona la prudencia de la revision: aumenta el peso de comprobar calendario, liquidez y confirmacion posterior antes de extraer conclusiones. Con riesgo macro marcado como alto, la lectura debe tratarse como sensible a evento y no como una validacion aislada del grafico.

\section*{Conclusion operativa y financiera}
Conclusion operativa y financiera: US30 queda como rsi extreme with context de revision, no como instruccion de mercado. La prioridad de revision indicada por el analisis es 4/5 y el riesgo macro figura como alto. La interpretacion conjunta debe apoyarse en tres filtros humanos: que el precio confirme el comportamiento descrito en el grafico, que las contradicciones de tendencia o estructura no invaliden el caso, y que no exista un evento macro cercano capaz de distorsionar el movimiento. Sin esas comprobaciones, el informe solo justifica seguimiento y documentacion, no ejecucion.

\section*{Comprobaciones humanas}
\begin{itemize}
  \item Verificar si US500 y US100 confirman el mismo deterioro o si US30 muestra debilidad relativa aislada.
  \item Comprobar una recuperacion real del precio sobre 51.116-51.237 antes de interpretar la sobreventa como estabilizacion tecnica.
  \item Revisar si el RSI H1 sale de sobreventa con mejora de estructura en M15, no solo por una pausa lateral.
  \item Contrastar el nivel 50.727-50.740 como referencia inferior cercana del rango/Fibonacci y minimo previo.
  \item Revisar calendario y titulares macro antes de cualquier evaluacion posterior, especialmente por el impacto del informe laboral y expectativas de tipos.
\end{itemize}

\section*{Fuentes}
\textbf{Fuentes locales}\par
\begin{tabularx}{\textwidth}{p{4.2cm}X}
setup\_context.json & Datos estructurados del setup \\
market\_context.json & Contexto agregado de mercado \\
ohlc\_window.csv & Ventana de velas OHLC \\
chart\_layers.csv & Capas dibujadas en el grafico \\
source\_manifest.json & Trazabilidad del paquete \\
chart.png & Imagen renderizada del grafico \\
\end{tabularx}

\vspace{0.45em}
\textbf{Enlaces externos}\par
\begin{tabularx}{\textwidth}{p{7.4cm}X}
\href{\detokenize{https://apnews.com/article/b9d2661cbba6cc326c618c06769d8291}}{https://apnews.com/article/b9d2661cbba6cc326c618c06769d8291} & Fuente externa consultada para contexto macro/noticias. \\
\href{\detokenize{https://www.investing.com/economic-calendar/nonfarm-payrolls-227}}{https://www.investing.com/economic-calendar/nonfarm-payrolls-227} & Fuente externa consultada para contexto macro/noticias. \\
\href{\detokenize{https://www.federalreserve.gov/monetarypolicy/fomccalendars.htm}}{https://www.federalreserve.gov/monetarypolicy/fomccalendars.htm} & Fuente externa consultada para contexto macro/noticias. \\
\href{\detokenize{https://www.axios.com/2026/06/05/interest-rates-warsh-inflation}}{https://www.axios.com/2026/06/05/interest-rates-warsh-inflation} & Fuente externa consultada para contexto macro/noticias. \\
\end{tabularx}

\end{document}
